\documentclass[11pt]{amsart}
\usepackage[english]{babel}
\usepackage[utf8]{inputenc}
\usepackage[T1]{fontenc}
\usepackage{amsmath, amssymb, amsfonts}
\usepackage{amsthm}
\usepackage{mathrsfs}
\usepackage{enumitem}
\usepackage{mathtools}
\usepackage{todonotes}
\usepackage{standalone}
\usepackage{pgfplots}
\pgfplotsset{compat=1.17}
\usepackage{thmtools, thm-restate}
\usepackage{booktabs}
\usepackage{hyperref}
\hypersetup{hidelinks}
\usepackage[capitalise, noabbrev]{cleveref}
\usepackage[font=small]{caption}
\usepackage{microtype}
\sloppy
% Preamble for LaTeX documents

% \hypersetup{
%   colorlinks   = true,
%   urlcolor     = darkred,
%   linkcolor    = darkred,
%   citecolor   = darkgreen}


\DeclarePairedDelimiter\ceil{\lceil}{\rceil}
\DeclarePairedDelimiter\floor{\lfloor}{\rfloor}
\DeclarePairedDelimiter\abs{\lvert}{\rvert}
\newcommand{\norm}[1]{\left\lVert#1\right\rVert}

\usepackage{bm}           % Bold math symbols
\usepackage{soul}
% Graphics packages (uncomment as needed)
\usepackage{graphicx}
\graphicspath{{./figures/}}
%\usepackage{caption}
% \usepackage{subcaption}
\usepackage{tikz}
\usepackage{tikz-cd}

% algorithm2e package for pseudocode
\usepackage[vlined, ruled,linesnumbered]{algorithm2e}
\usepackage{algpseudocode}

% cleveref package (for enhanced cross-referencing)
%\usepackage[noabbrev,capitalise]{cleveref}

% Hyperlinks and colors (uncomment as needed)
%\usepackage[colorlinks=true, linkcolor=blue, filecolor=magenta, urlcolor=cyan]{hyperref}

% Custom commands (add your own as needed)
%\newcommand{\mycommand}{...}


% \newtheorem{theorem}{Theorem}[section]
% \newtheorem{lemma}[theorem]{Lemma}
% \newtheorem{proposition}[theorem]{Proposition}
% \newtheorem{corollary}[theorem]{Corollary}
% \newtheorem*{corollary*}{Corollary}
% \newtheorem{conjecture}[theorem]{Conjecture}
% \newtheorem{atheorem}{Theorem}
% \renewcommand{\theatheorem}{\Alph{atheorem}}
% \newtheorem{acorollary}[atheorem]{Corollary}
% \renewcommand{\theacorollary}{\Alph{acorollary}}
% \newtheorem{question}[theorem]{Question}
% \theoremstyle{definition}
% \newtheorem{definition}[theorem]{Definition}
% \newtheorem{notation}[theorem]{Notation}
% \newtheorem{convention}[theorem]{Convention}
% \newtheorem{example}[theorem]{Example}
% \newtheorem{examples}[theorem]{Examples}
% \theoremstyle{remark}
% \newtheorem{remark}[theorem]{Remark}
% \newtheorem{claim}[theorem]{Claim}
% \newtheorem{discussion}[theorem]{Discussion}
% \newtheorem{computation}[theorem]{Computation}


% Notation for math sets
\newcommand{\N}{\mathbb{N}} % Natural numbers
\newcommand{\Z}{\mathbb{Z}} % Integers
\newcommand{\Q}{\mathbb{Q}} % Rational numbers
\newcommand{\R}{\mathbb{R}} % Real numbers
\newcommand{\C}{\mathbb{C}} % Complex numbers
\newcommand{\mat}[1]{\mathbf{#1}} % Matrix notation (bold)

\newcommand{\define}[1]{\emph{#1}}

% using \varepsilon as default
\let\uglyepsilon\epsilon
\let\epsilon\varepsilon

\newcommand{\Ball}{\text{Ball}}
\DeclareMathOperator{\Nrv}{Nrv}
\DeclareMathOperator{\V}{\mathcal{V}}
\DeclareMathOperator{\Vor}{Vor}
\DeclareMathOperator{\DelT}{Del}
\newcommand{\DC}{\mathcal{DC}}
\newcommand{\NN}{\mathcal{N}}
\DeclareMathOperator{\Del}{\mathcal{D}}
\DeclareMathOperator{\Tel}{Tel}
\DeclareMathOperator{\TelVor}{\mathcal{T}}
\DeclareMathOperator{\TelBall}{\mathcal{T}B}
\DeclareMathOperator{\TelOffset}{\mathcal{TO}}
\DeclareMathOperator{\TelCech}{\mathcal{TC}}
\DeclareMathOperator{\TelDel}{\mathcal{TD}}
\DeclareMathOperator{\TelDelCech}{\mathcal{TDC}}
\DeclareMathOperator{\TelVO}{\mathcal U}
\DeclareMathOperator{\TelDC}{\mathcal W}

\DeclareMathOperator{\Cech}{\mathcal{C}}
\DeclareMathOperator{\DelCech}{\mathcal{DC}}
\DeclareMathOperator{\Grid}{\mathrm{Grid}}
\DeclareMathOperator{\Star}{Star}
\DeclareMathOperator{\Lk}{link}
\DeclareMathOperator{\im}{im}
\DeclareMathOperator{\colim}{colim}

\DeclareMathOperator*{\minimize}{minimize}



% category notations
\newcommand{\Int}{\textbf{Int}}
\newcommand{\Top}{\textbf{Top}}
\newcommand{\Simp}{\textbf{Simp}}

\newcommand{\cover}{\textnormal{Cov}}
\newcommand{\filt}{{\bullet}}



\PassOptionsToPackage{x11names}{xcolor}
\usepackage{xcolor}
\definecolor{darkred}{RGB}{139,0,0}
\definecolor{darkblue}{RGB}{0,0,139}
\definecolor{darkgreen}{RGB}{0,100,0}

\newcommand{\deff}{\emph}
\newcommand{\Delaunay}[0]{\mathcal{D}}
\newcommand{\Offset}[0]{\mathcal{O}}
\newcommand{\Conv}{\mathrm{Conv}}
\newcommand{\Incr}[0]{\mathcal{I}}
\newcommand{\DegDel}[0]{\mathcal{DD}}
\newcommand{\LiftVor}{\mathcal{V}}
\newcommand{\degpar}[0]{\Delta}
\newcommand{\St}{\mathrm{Star}}

\providecommand\given{}
\newcommand\SetSymbol[1][]{%
  \nonscript\:#1\vert
  \allowbreak
  \nonscript\:
  \mathopen{}}
\DeclarePairedDelimiterX\Set[1]\{\}{%
  \renewcommand\given{\SetSymbol[]}
  #1
}

\makeatletter
\DeclareRobustCommand{\subto}{%
  \mathrel{\mathpalette\short@to\relax}%
}

\newcommand{\short@to}[2]{%
  \mkern2mu
  \clipbox{{.35\width} 0 0 0}{$\m@th#1\vphantom{+}{\rightarrow}$}%
  }
\makeatother

\makeatletter
\DeclareRobustCommand{\subleft}{%
  \mathrel{\mathpalette\short@left\relax}%
}

\newcommand{\short@left}[2]{%
  \mkern2mu
  \clipbox{0 0 {.35\width} 0}{$\m@th#1\vphantom{+}{\leftarrow}$}%
  }
\makeatother

\makeatletter
\DeclareRobustCommand{\subup}{%
  \mathrel{\mathpalette\short@up\relax}%
}

\newcommand{\short@up}[2]{%
  \mkern2mu
  \clipbox{0 {.35\height} 0 0}{$\m@th#1\vphantom{+}{\uparrow}$}%
}
\makeatother

\makeatletter
\DeclareRobustCommand{\subdown}{%
  \mathrel{\mathpalette\short@down\relax}%
}
\newcommand{\short@down}[2]{%
  \mkern2mu
  \clipbox{0 0 0 {.45\height}}{$\m@th#1\vphantom{+}{\downarrow}$}%
}
\makeatother

\makeatletter
\DeclareRobustCommand{\subupright}{%
  \mathrel{\mathpalette\short@near\relax}%
}
\newcommand{\short@near}[2]{%
  \mkern2mu
  \clipbox{{.55\width} {.55\height} 0 0}{$\m@th#1\vphantom{+}{\nearrow}$}%
}
\makeatother

\makeatletter
\DeclareRobustCommand{\subdownleft}{%
  \mathrel{\mathpalette\short@sw\relax}%
}
\newcommand{\short@sw}[2]{%
  \mkern2mu
  \clipbox{0 0 {.55\width} {.55\height}}{$\m@th#1\vphantom{+}{\swarrow}$}%
}
\makeatother

\newcommand{\angel}[1]{\textcolor{darkblue}{#1}}
\newcommand{\tung}[1]{\textcolor{purple}{#1}}
\newcommand{\abhishek}[1]{\textcolor{darkgreen}{#1}}

\renewcommand{\subparagraph}[1]{\subsubsection*{\textbf{\textup{#1}}}}

\theoremstyle{plain}
\newtheorem{theorem}{Theorem}[section]
\newtheorem{lemma}[theorem]{Lemma}
\newtheorem{corollary}[theorem]{Corollary}
\newtheorem{proposition}[theorem]{Proposition}
\newtheorem{claim}[theorem]{Claim}

\theoremstyle{definition}
\newtheorem{definition}[theorem]{Definition}
\newtheorem{example}[theorem]{Example}

\theoremstyle{remark}
\newtheorem{remark}[theorem]{Remark}
\newtheorem{observation}[theorem]{Observation}

%\theoremstyle{plain}
%\newtheorem{theorem}{Theorem}
%\newtheorem{lemma}[theorem]{Lemma}
%\newtheorem{corollary}[theorem]{Corollary}
%\newtheorem{proposition}[theorem]{Proposition}
%\newtheorem{claim}[theorem]{Claim}
%
%\theoremstyle{definition}
%\newtheorem{definition}[theorem]{Definition}
%\newtheorem{example}[theorem]{Example}
%
%\theoremstyle{remark}
%\newtheorem{remark}[theorem]{Remark}
%\newtheorem{observation}[theorem]{Observation}

\newcommand{\SizeBound}{{\lceil (d+1)/2 \rceil + 1}}
\newcommand{\CompBound}{{\lceil d/2\rceil +2}}

\usepackage{trimclip}
\usepackage{fontspec}
\setmainfont{DejaVu Serif}
\setlength{\emergencystretch}{3em}
